% Part: model-theory
% Chapter: interpolation
% Section: definability

\documentclass[../../../include/open-logic-section]{subfiles}

\begin{document}

\olfileid[hi]{mod}{int}{def}

\olsection{परिभाषणीयता प्रमेय}

अंतर्वेशन प्रमेय का एक महत्त्वपूर्ण अनुप्रयोग बेथ का परिभाषणीयता प्रमेय
है। किसी $n$-स्थानी संबंध~$R$ को परिभाषित करने के लिए हम ऐसा
!!a{formula}~$!C$ दे सकते हैं जिसमें $n$ मुक्त !!{variable}s हों और
जिसमें~$R$ न आता हो। यह~$R$ की,~$!C$ के पदों में, एक \emph{स्पष्ट}
परिभाषा होगी। तब हम यह भी कह सकते हैं कि कोई सिद्धांत~$\Sigma(P)$,
किसी !!{language} में जिसमें $n$-स्थानी !!{predicate}~$P$ हो,~$P$ को
स्पष्टतः परिभाषित करता है, यदि उसमें कोई
औपचारिक स्पष्ट परिभाषा हो (या कम-से-कम वह उससे तात्पर्यित हो), अर्थात्
\[
\Sigma(P) \Entails \lforall[x_1][\dots
  \lforall[x_n][(\Atom{P}{x_1,\dots, x_n} \liff !C(x_1, \dots,
    x_n))]].
\]
किंतु स्पष्ट परिभाषा किसी संबंध को परिभाषित करने---अर्थात् उसे पूरी तरह
निर्धारित करने---का केवल एक ढंग है। कोई सिद्धांत ऐसा भी हो सकता है कि
किसी भी प्रतिरूप में शेष !!{language} की व्याख्या,~$P$ की व्याख्या को
निश्चित कर दे। परिभाषणीयता प्रमेय कहता है कि जब भी कोई सिद्धांत इस प्रकार
~$P$ की व्याख्या निश्चित करता है---जब भी वह~$P$ को \emph{अंतर्निहित रूप
से परिभाषित करता है}---तब वह उसे स्पष्टतः भी परिभाषित करता है।

\begin{defn}
मान लें कि $\Lang{L}$ ऐसी !!a{language} है जिसमें !!{predicate}~$P$ नहीं
है। कोई समुच्चय $\Sigma(P)$, $\Lang{L}
\cup \{P\}$ के !!{sentence}s का,~$P$ को तभी और केवल तभी
\emph{स्पष्टतः परिभाषित करता है} जब ऐसा !!a{formula}
~$!C(x_1, \dots, x_n)$, $\Lang{L}$ का, हो कि
\[
\Sigma(P) \Entails \lforall[x_1][\dots
  \lforall[x_n][(\Atom{P}{x_1,\dots, x_n} \liff !C(x_1, \dots,
    x_n))]].
\]
\end{defn}

\begin{defn}
मान लें कि $\Lang{L}$ ऐसी !!a{language} है जिसमें !!{predicate}s~$P$
और~$P'$ नहीं हैं। कोई समुच्चय $\Sigma(P)$, $\Lang{L} \cup \{P\}$ के
!!{sentence}s का, $P$ को तभी और केवल तभी \emph{अंतर्निहित रूप से
परिभाषित करता है} जब
\[
\Sigma(P) \cup \Sigma(P') \Entails \lforall[x_1][\dots
  \lforall[x_n][(\Atom{P}{x_1,\dots, x_n} \liff \Atom{P'}{x_1,\dots,
      x_n})]],
\]
जहाँ $\Sigma(P')$ उस प्रतिस्थापन का परिणाम है जिसमें $P$ को $P'$ से
$\Sigma(P)$ में सर्वत्र बदला जाता है।
\end{defn}

दूसरे शब्दों में, किसी भी प्रतिरूप $\Struct{M}$ और $R, R' \subseteq
\Domain{M}^n$ के लिए, यदि $\Expan{M}{R} \Entails \Sigma(P)$ और
$\Expan{M}{R'} \Entails \Sigma(P')$ दोनों हों, तो $R=R'$; यहाँ
$\Expan{M}{R}$ वह !!{structure}~$\Struct{M'}$ है जो $\Lang{L}$ के
$\Lang{L} \cup \{P\}$ तक विस्तार के लिए ऐसी है कि
$\Assign{P}{M'} = R$, और
$\Expan{M}{R'}$ भी इसी प्रकार है।

\begin{thm}[बेथ का परिभाषणीयता प्रमेय] कोई समुच्चय $\Sigma(P)$,
  $\Lang{L}
  \cup\{P\}$-!!{formula}s का, $P$ को अंतर्निहित रूप से तभी और केवल तभी
  परिभाषित करता है जब $\Sigma(P)$, $P$ को स्पष्टतः परिभाषित करता है।
\end{thm}

\begin{proof}
यदि $\Sigma(P)$, $P$ को स्पष्टतः परिभाषित करता है, तो दोनों
\begin{align*}
  \Sigma(P) & \Entails & \lforall[x_1][\dots \lforall[x_n]
    [(\Atom{P}{x_1,\dots, x_n} \liff !C(x_1,\dots,x_n))]]\\
  \Sigma(P') & \Entails & \lforall[x_1][\dots \lforall[x_n]
    [(\Atom{P'}{x_1,\dots, x_n} \liff !C(x_1,\dots,x_n))]]
\end{align*}
और निष्कर्ष निकलता है। विलोम के लिए मान लें कि $\Sigma(P)$, $P$ को
अंतर्निहित रूप से परिभाषित करता है। पहले हम !!{constant}s $c_1$,
\dots,~$c_n$ को $\Lang{L}$ में जोड़ते हैं। तब
\[
\Sigma(P) \cup \Sigma(P') \Entails
\Atom{P}{c_1, \dots, c_n} \to  \Atom{P'}{c_1, \dots, c_n}.
\]
संहतता से ऐसे परिमित समुच्चय $\Delta_0 \subseteq \Sigma(P)$ और
$\Delta_1 \subseteq \Sigma(P')$ हैं कि
\[
\Delta_0 \cup \Delta_1 \Entails
\Atom{P}{c_1, \dots, c_n} \to \Atom{P'}{c_1, \dots, c_n}.
\]
$!D(P)$ को उन सभी !!{sentence}s $!A(P)$ का संयोजन मानें जिनके लिए या तो
$!A(P) \in \Delta_0$ या $!A(P') \in \Delta_1$ है; और $!D(P')$ को उन
सभी !!{sentence}s $!A(P')$ का संयोजन मानें जिनके लिए या तो
$!A(P) \in \Delta_0$ या $!A(P') \in \Delta_1$ है। तब $!D(P) \land
!D(P') \Entails \Atom{P}{c_1, \dots, c_n} \to P'c_1\dots c_n$। इसे
हम इस प्रकार पुनर्व्यवस्थित कर सकते हैं कि प्रत्येक !!{predicate},
$\Entails$ के एक ही ओर आए:
\[
!D(P) \land \Atom{P}{c_1, \dots, c_n} \Entails
!D(P') \to \Atom{P'}{c_1, \dots, c_n}.
\]
क्रेग के अंतर्वेशन प्रमेय से ऐसा !!a{sentence} $!C(c_1,\dots, c_n)$ है
जिसमें $P$ या $P'$ नहीं आता और जिसके लिए:
\begin{align*}
  {}!D(P) \land \Atom{P}{c_1, \dots, c_n} & \Entails !C(c_1,\dots, c_n); \\
  {}!C(c_1,\dots, c_n) & \Entails !D(P') \to \Atom{P'}{c_1, \dots, c_n}.
\end{align*}
इन दोनों तात्पर्यों में पहले से हमें मिलता है: $!D(P) \Entails
\Atom{P}{c_1,\dots, c_n} \lif !C(c_1,\dots, c_n)$। और दूसरे से, चूँकि
कोई $\Lang{L} \cup \{P\}$-प्रतिरूप $\Expan{M}{R}
\Entails !A(P)$ तभी और केवल तभी जब संगत $\Lang{L} \cup
\{P'\}$-प्रतिरूप $\Expan{M}{R} \models !A(P')$, हमें
$!C(c_1,\dots, c_n) \Entails !D(P) \lif \Atom{P}{c_1,\dots, c_n}$
मिलता है, जिससे:
\[
!D(P) \Entails !C(c_1,\dots,c_n) \to \Atom{P}{c_1,\dots, c_n}.
\]
दोनों को मिलाने पर $!D(P) \Entails \Atom{P}{c_1,\dots, c_n}
\liff !C(c_1, \dots, c_n)$, और एकदिशता तथा व्यापकीकरण द्वारा यह भी:
\[
\Sigma(P) \Entails
\lforall[x_1][\dots\lforall[x_n][(\Atom{P}{x_1,\dots, x_n} \liff
    !C(x_1,\dots, x_n))]].
\]
\end{proof}

\end{document}
